\documentclass[aspectratio=169,11pt]{beamer}
% Shared CMPSC 480 Beamer preamble.
\usepackage[T1]{fontenc}
\usepackage[utf8]{inputenc}
\usepackage{lmodern}
\usepackage{amsmath,amssymb,mathtools}
\usepackage{booktabs,array,tabularx}
\usepackage{xcolor}
\usepackage{listings}
\usepackage{tikz}
\usetikzlibrary{arrows.meta,positioning,shapes.geometric}

% Ignore only negligible TeX box diagnostics that are below visual tolerance.
% Larger layout defects still appear in the build log and in rendered QA.
\vfuzz=3pt
\hbadness=2000

\definecolor{courseink}{HTML}{111111}
\definecolor{coursegray}{HTML}{EDEDED}
\definecolor{courserule}{HTML}{B8BCC4}
\definecolor{courseblue}{HTML}{3D8DFF}
\definecolor{courselightblue}{HTML}{D0EDFA}
\definecolor{coursemuted}{HTML}{5C6470}
\definecolor{coursegreen}{HTML}{0B7A53}
\definecolor{coursered}{HTML}{B3261E}

\usetheme{default}
\usefonttheme{professionalfonts}
\setbeamertemplate{navigation symbols}{}
\setbeamersize{text margin left=0.65cm,text margin right=0.65cm}

\setbeamercolor{normal text}{fg=courseink,bg=white}
\setbeamercolor{frametitle}{fg=courseink,bg=white}
\setbeamercolor{title}{fg=courseink,bg=white}
\setbeamercolor{subtitle}{fg=coursemuted,bg=white}
\setbeamercolor{structure}{fg=courseblue}
\setbeamercolor{alerted text}{fg=coursered}
\setbeamercolor{block title}{fg=courseink,bg=courselightblue}
\setbeamercolor{block body}{fg=courseink,bg=coursegray}
\setbeamercolor{example text}{fg=coursegreen}

\setbeamerfont{title}{size=\fontsize{25}{29}\selectfont,series=\bfseries}
\setbeamerfont{subtitle}{size=\fontsize{13}{16}\selectfont}
\setbeamerfont{frametitle}{size=\fontsize{19}{22}\selectfont,series=\bfseries}
\setbeamerfont{normal text}{size=\fontsize{11}{14}\selectfont}
\setbeamerfont{footline}{size=\fontsize{7.5}{9}\selectfont}

\setbeamertemplate{frametitle}{%
  \nointerlineskip
  % Use content-driven height so a long assessment title can wrap without
  % being clipped by a fixed-height title bar.
  \begin{beamercolorbox}[wd=\paperwidth,sep=0.17cm,leftskip=0.48cm,rightskip=0.48cm]{frametitle}
    \insertframetitle
  \end{beamercolorbox}
  \vspace{-0.08cm}
  {\color{courserule}\rule{\textwidth}{0.35pt}}
}

\setbeamertemplate{footline}{%
  \leavevmode
  \begin{beamercolorbox}[wd=\paperwidth,ht=2.6ex,dp=1.1ex,leftskip=.65cm,rightskip=.65cm]{author in head/foot}
    \color{coursemuted}\insertshorttitle\hfill\insertframenumber/\inserttotalframenumber
  \end{beamercolorbox}
}

\setbeamertemplate{itemize item}{\color{courseblue}\small$\blacktriangleright$}
\setbeamertemplate{itemize subitem}{\color{courseblue}\scriptsize$\bullet$}
\setbeamertemplate{enumerate item}{\color{courseblue}\insertenumlabel.}

\lstdefinestyle{coursepython}{
  language=Python,
  basicstyle=\ttfamily\fontsize{8.5}{10}\selectfont,
  keywordstyle=\color{courseblue}\bfseries,
  commentstyle=\color{coursemuted}\itshape,
  stringstyle=\color{coursegreen},
  showstringspaces=false,
  columns=fullflexible,
  keepspaces=true,
  frame=single,
  rulecolor=\color{courserule},
  backgroundcolor=\color{coursegray},
  breaklines=true,
  tabsize=4,
  xleftmargin=0.15cm,
  xrightmargin=0.15cm,
  framexleftmargin=0.15cm,
  framexrightmargin=0.15cm
}
\lstset{style=coursepython}

\newcommand{\points}[1]{\hfill{\color{courseblue}\textbf{[#1 points]}}}
\newcommand{\deliverable}[1]{\textbf{\color{courseblue}#1}}
\newcommand{\code}[1]{\texttt{#1}}
\newcommand{\keyidea}[1]{%
  \begin{beamercolorbox}[rounded=false,sep=0.55em,wd=\linewidth]{block body}
    \textbf{Key idea.} #1
  \end{beamercolorbox}
}
\newcommand{\answerlabel}{\textbf{\color{coursegreen}Answer.}\ }
\newcommand{\gradinglabel}{\textbf{\color{courseblue}Grading guidance.}\ }

\newenvironment{questionblock}[1]
  {\begin{block}{#1}}
  {\end{block}}

\newenvironment{solutionblock}[1]
  {\begin{block}{#1}}
  {\end{block}}

\AtBeginSection[]{
  \begin{frame}
    \vfill
    {\usebeamerfont{title}\color{courseink}\insertsectionhead\par}
    \vspace{0.4cm}
    {\color{courseblue}\rule{0.32\paperwidth}{3pt}}
    \vfill
  \end{frame}
}


% CMPSC480_TOTAL_POINTS: 10
% CMPSC480_ITEM_IDS: 1,2,3,4
\title[Week 5 Assignment Solutions]{Week 5 Assignment: Tabular Q-Learning}
\subtitle{A tested model-free agent with controlled exploration and oracle comparison}
\author{CMPSC 480 - Instructor Solution Deck}
\date{}

\begin{document}


\begin{frame}
  \titlepage
\vspace{-0.35cm}
\begin{center}
\color{coursered}\bfseries Instructor material - do not distribute with the question deck
\end{center}
\end{frame}

\begin{frame}{Solution overview}
\begin{columns}[T,onlytextwidth]
\begin{column}{0.62\textwidth}
\Large This instructor deck provides a complete matched solution and grading guidance.

\vspace{0.35cm}
\normalsize
Coverage: Week 5: TD learning, exploration, terminal targets, and empirical evaluation
\end{column}
\begin{column}{0.33\textwidth}
\begin{block}{Points}10 total\end{block}
\begin{block}{Expected time}6-8 hours\end{block}
\begin{block}{Submission}Repository, learning curves, and report\end{block}
\end{column}
\end{columns}
\end{frame}

\begin{frame}{Instructor verification checklist}
\small
\begin{itemize}
  \item Use Python 3.11 and NumPy only unless the prompt explicitly permits another package.
  \item Preserve the published interfaces and make all tie-breaking deterministic.
  \item State assumptions, validate inputs, and handle empty, terminal, and unreachable cases.
  \item Include focused tests whose names explain the defect each test would expose.
  \item Report measured evidence; do not replace experimental results with unsupported claims.
  \item Seed every source of randomness used in reported experiments.
\end{itemize}
\end{frame}

\begin{frame}{Solution 1.a - Agent and exploration contract (1 points)}
\small
\textbf{Question focus.} Define observe, select\_action, update, and evaluate-mode behavior.

\vspace{0.25cm}
\answerlabel The agent receives transitions from the environment, stores Q values keyed by state-action, selects only legal actions, updates from observations, and uses epsilon zero during evaluation.
\end{frame}

\begin{frame}{Solution 1.b - Agent and exploration contract (1 points)}
\small
\textbf{Question focus.} Specify epsilon-greedy selection and a bounded decay schedule.

\vspace{0.25cm}
\answerlabel With probability epsilon choose uniformly among legal actions; otherwise choose a stable argmax. Decay is applied at a documented cadence and clipped to epsilon\_min.
\end{frame}

\begin{frame}{Solution 2.a - Q-learning update (1 points)}
\small
\textbf{Question focus.} State the update and define the TD error.

\vspace{0.25cm}
\answerlabel The target is r plus gamma times max over next legal actions Q(s-prime,a-prime) for nonterminal transitions. Delta is target minus Q(s,a), and Q(s,a) increases by alpha times delta.
\end{frame}

\begin{frame}{Solution 2.b - Q-learning update (1 points)}
\small
\textbf{Question focus.} Compute one update for Q=2, r=1, gamma=0.9, max next Q=4, and alpha=0.5.

\vspace{0.25cm}
\answerlabel The target is 1+0.9*4=4.6; delta is 2.6; the updated value is 2+0.5*2.6=3.3.
\end{frame}

\begin{frame}{Solution 2.c - Q-learning update (1 points)}
\small
\textbf{Question focus.} Explain and test terminal handling.

\vspace{0.25cm}
\answerlabel For done=True the target is r only; no next-state maximum is used. A test gives a very large next Q value and confirms it does not change the terminal update.
\end{frame}

\begin{frame}{Solution 3.a - Training protocol (1 points)}
\small
\textbf{Question focus.} Describe reset, interaction, update, and episode termination order.

\vspace{0.25cm}
\answerlabel Reset returns the initial state; each step selects, executes, records reward and done, updates once, and advances state. The loop ends on done or a maximum-step truncation.
\end{frame}

\begin{frame}{Solution 3.b - Training protocol (1 points)}
\small
\textbf{Question focus.} Log return, episode length, epsilon, and moving-average return.

\vspace{0.25cm}
\answerlabel Metrics use consistent windows and distinguish truncation from environment termination. Seeds and hyperparameters are recorded with the run.
\end{frame}

\begin{frame}{Solution 4.a - Evaluation and tests (1 points)}
\small
\textbf{Question focus.} Evaluate the greedy policy without learning or exploration.

\vspace{0.25cm}
\answerlabel Set epsilon to zero, freeze Q values, run independent seeded episodes, and report mean return, success rate, and confidence or variability.
\end{frame}

\begin{frame}{Solution 4.b - Evaluation and tests (1 points)}
\small
\textbf{Question focus.} Compare greedy actions or returns with a value-iteration oracle.

\vspace{0.25cm}
\answerlabel Map both policies to the same MDP semantics. Count exact or tie-equivalent action agreement and explain differences caused by incomplete visitation or stochastic estimates.
\end{frame}

\begin{frame}{Solution 4.c - Evaluation and tests (1 points)}
\small
\textbf{Question focus.} Run one controlled hyperparameter ablation and cover sparse-reward and unseen-state cases.

\vspace{0.25cm}
\answerlabel Vary only alpha, epsilon schedule, or gamma. Report curves and explain bias/noise effects. Unseen states start from a documented default; sparse rewards do not cause crashes or false convergence claims.
\end{frame}

\begin{frame}{Edge-case reasoning and expected behavior}
\small
\begin{itemize}
  \item Treat it as terminal or raise a domain-contract error.
  \item Use target equal to the immediate reward.
  \item Choose deterministically in evaluation and uniformly only in exploration.
  \item Use the declared initialization value for every legal action.
  \item Record truncation separately and stop without a synthetic terminal bootstrap.
\end{itemize}
\end{frame}

\begin{frame}{Grading rubric}
\scriptsize
\begin{tabularx}{\textwidth}{>{\bfseries}p{0.28\textwidth} p{0.08\textwidth} X}
\toprule
Category & Points & Full-credit evidence \\
\midrule
Agent and exploration contract & 2 & Separation, legal actions, exploration probabilities, and ties are correct. \\
Q-learning update & 3 & TD target, terminal mask, and in-place assignment are correct. \\
Training protocol & 2 & Episodes terminate safely and training/evaluation are distinct. \\
Evaluation and tests & 3 & Metrics are unbiased, controlled, and interpreted correctly. \\
\bottomrule
\end{tabularx}
\end{frame}

\begin{frame}{Reflection exemplar}
\small
\answerlabel A strong response verifies shared MDP semantics first, then uses visitation counts, TD error, or repeated-run evidence to attribute the remaining gap to exploration, finite data, or hyperparameters.

\vspace{0.25cm}
\gradinglabel Award credit for a specific claim supported by technical evidence from the submitted work.
\end{frame}

\end{document}
